\chapter{AI Hardware Architecture and Compiler Constraints}
\label{chap:ch03}
\typeout{START_CHAPTER "ch03" \theabspage}

\section{Why Lab-Tuned Kernels Fail in Production}
\label{sec:industrial_hook}

Many compiler teams have shipped fusion kernels that hit 90\%+ utilization in the lab, then watched production latency double after deploy.
The root cause is often not the math---it is a few kilobytes of shared memory over budget that silently spills activations to HBM, wiping out fusion gains \citep{memoryfloor2026,nvidia2022hopper}.
Chapter~2 taught dataflow-first design; \textbf{every dataflow plan must land on concrete hardware constraints}---fusion without byte budgets is an architecture fiction.
``The GPU has fast memory'' is useless; ``an H100 SM exposes up to 228\,KB shared memory and a 256\,KB register file'' is input a tiling pass can solve against.

\paragraph{Bridge from Chapter~2.}
Dataflow mindset names \emph{where} tensors should live; this chapter names \emph{what fits} on each silicon class.
A fusion region that needs 250\,KB on-chip cannot exist on an H100 thread block without spilling---the schedule is wrong before anyone tunes MMA tiles \citep{yirage2026,chen2020compilerbasedhardw}.
Decode optimization is never ``roughly aware'' of hardware; a few kilobytes separate a 2$\times$ win from a regression \citep{taxbreak2026}.

\paragraph{Factory analogy (memory hierarchy).}
Registers are parts in the worker's hands---instant access, tiny capacity.
Shared memory / on-chip SRAM is the workbench beside the station.
L2 / LLC is the shop-floor storeroom a short walk away.
HBM / DDR is the remote warehouse across the plant \citep{dlbook,memoryfloor2026}.
Decode wins come from keeping activations off the warehouse route; hardware constraints tell you how much can stay on the bench.

\paragraph{Composite tri-hardware case (7B dense layer).}
The same mathematics, three deployment outcomes (composite internal migrations):
\textbf{GPU}---operator-first 42\,ms/token at batch-1; hardware-aware fusion ${\sim}18$\,ms/token (${\sim}2.3$\times$) when shared-memory tiling respects SM caps \citep{memoryfloor2026,taxbreak2026}.
\textbf{CPU}---naively porting the GPU fusion evicted KV from LLC and ran 15\% \emph{slower} than native blocking; cache-aware retile gained ${\sim}35\%$ over eager PyTorch \citep{patel2025llminference,rotem2018glow}.
\textbf{NPU}---dynamic per-op dispatch did not legalize; static ADF tile+DMA plan compiled and beat a same-power GPU class by ${\sim}40\%$ on per-token latency in one edge SKU \citep{amd2024xdna,amd2024aie}.
Matching architecture yields gains; ignoring constraints yields disasters \citep{yirage2026}.

\paragraph{What this chapter gives you.}
You will leave able to (1)~read a SKU datasheet row as a compiler constraint, (2)~predict failure modes when a fusion violates capacity, (3)~pick the right debug tool per class, and (4)~justify target-specific schedules with bytes/token and launches/token instead of peak FLOPs \citep{taxbreak2026,memoryfloor2026,patel2025llminference}.
Chapter~2 told you to plan residency; this chapter tells you which bench has how many kilobytes \citep{yirage2026,nvidia2022hopper,amd2024xdna}.

\paragraph{Reader prerequisites.}
We assume Chapter~1's goodput counters and Chapter~2's export-edge mindset \citep{taxbreak2026,memoryfloor2026,chen2020compilerbasedhardw}.
No chip-design background required---but you must be willing to read capacities as hard inequalities, not inspirational slogans \citep{patel2025llminference}.
If SIMT or NUMA are new, the terminology bridge in \S\ref{sec:hardware_landscape} defines them on first use \citep{nvidia2022hopper,dlbook}.

\paragraph{How to read this chapter.}
Skim Table~\ref{tab:tri_hw_decode_compare} first if you are time-constrained---it is the cheat sheet the rest of the chapter elaborates \citep{patel2025llminference}.
GPU engineers should still read \S\ref{sec:npu_constraints} to understand why fleet schedules diverge; edge engineers should read \S\ref{sec:gpu_constraints} to see what flexibility they traded away \citep{amd2024xdna,nvidia2022hopper}.
Software-first readers: treat every capacity number as an inequality in the compiler, not trivia for a hardware quiz \citep{yirage2026,chen2020compilerbasedhardw}.

\paragraph{Skillful engineering over SKU shopping.}
Once frameworks and famous kernels are enabled, the remaining lever is schedules that fit silicon \citep{yirage2026,memoryfloor2026}.
A mid-tier GPU with legal fusion beats a flagship running an eager loop---Memory-Floor's L4 versus H100 cells demonstrate the pattern \citep{memoryfloor2026,taxbreak2026}.
Hardware literacy is how you reproduce that pattern on \emph{your} fleet \citep{patel2025llminference}.

\paragraph{Spill signatures in telemetry.}
Production regressions often show flat GPU utilization with rising HBM throughput---classic spill or export-edge behavior \citep{memoryfloor2026,nvidia2022hopper}.
CPU regressions show rising LLC miss rate per token without code changes in the math kernel \citep{dlbook,patel2025llminference}.
NPU regressions show rising DDR bytes per token when a new ADF version added a DMA stage \citep{amd2024xdna,yirage2026}.
Learn to map telemetry shapes to constraint classes before opening profiler captures \citep{taxbreak2026,memoryfloor2026}.

\paragraph{Chapter artifact.}
By the end of this chapter you should maintain a one-page constraint sheet per deployment target: capacities, bottleneck column from Table~\ref{tab:tri_hw_decode_compare}, and baseline bytes/token plus launches/token \citep{patel2025llminference,taxbreak2026}.
Update the sheet when SKUs, context limits, or quantization policy change \citep{yirage2026,memoryfloor2026}.
Chapter~4 assumes that sheet exists and adds Hopper-specific knobs on top \citep{nvidia2022hopper,semianalysis2024tma}.

\paragraph{Closing reminder.}
Hardware constraints are not obstacles to clever math---they are the coordinate system clever schedules are written in \citep{chen2020compilerbasedhardw,yirage2026}.
Respect the numbers early; profilers should confirm wins, not discover basic illegality \citep{memoryfloor2026,patel2025llminference}.
Keep the constraint sheet in version control next to the compiler config \citep{yirage2026}.

\index{Hardware architecture}
\index{Compiler constraints}
\index{Memory hierarchy}

\section{AI Hardware Landscape: Three Architectures, Three Optimization Logics}
\label{sec:hardware_landscape}

% AUTO_VISUAL:tab_hw_landscape
\begin{table}[t]
\centering
\caption{AI accelerator classes and primary execution models.}
\label{tab:hw_landscape}
\begin{tabular}{lcc}
\toprule
\visTableHeader
\textbf{Class} & \textbf{Execution model} & \textbf{Decode fit} \\
\midrule
NVIDIA GPU & SIMT + tensor cores & Graphs, deep fusion \\
AMD GPU & Wavefront + matrix cores & HIP fusion \\
x86/ARM CPU & Cache + SIMD & Cache-blocked fusion \\
Edge NPU & Spatial dataflow & Static MegaKernel \\
\bottomrule
\end{tabular}
\end{table}

Production is no longer GPU-only.
An \textbf{NVIDIA GPU} is a \textbf{SIMT} (single-instruction multiple-thread) machine: thousands of threads scheduled dynamically over registers, shared memory, L2, and HBM \citep{nvidia2022hopper}.
An \textbf{AMD GPU} uses wavefronts and matrix cores via HIP/ROCm with similar hierarchy, different tooling \citep{amd2024aie}.
An \textbf{x86/ARM CPU} runs coherent caches and \textbf{SIMD} vector lanes with flexible control flow \citep{dlbook}.
An edge \textbf{NPU} such as AMD XDNA is a spatial dataflow fabric: fixed tiles, explicit SRAM, no hardware cache \citep{amd2024xdna}.

Each class fails batch-1 decode differently.
GPUs pay \textbf{kernel launch overhead} (scheduling overhead) per operator boundary---Chapter~1's launches/token \citep{taxbreak2026}.
CPUs trade peak bandwidth for flexible blocking and low idle cost \citep{patel2025llminference}.
NPUs trade flexibility for deterministic on-chip bandwidth when the static plan fits \citep{amd2024aie}.
Table~\ref{tab:hw_landscape} ``decode fit'' prescribes the artifact each class wants: captured fusion on GPU, cache-blocked bundles on CPU, static MegaKernel on NPU \citep{mpk2025,rotem2018glow}.

\paragraph{Terminology introduced here.}
\textbf{WGMMA} (warp-group matrix multiply-accumulate) is Hopper's high-throughput tensor-core instruction issued by warpgroups of four warps \citep{shah2024flashattention3,nvidia2022hopper}.
\textbf{DMA} (direct memory access) engines move data without burning compute threads---critical on NPUs and increasingly on GPUs via \textbf{TMA} \citep{semianalysis2024tma,amd2024xdna}.
\textbf{TMA} (Tensor Memory Accelerator) performs async copies between global memory and shared memory on Hopper+, overlapping movement with math \citep{nvidia2022hopper,semianalysis2024tma}.

\paragraph{Why the landscape fragmented.}
For a decade, vendor libraries (cuDNN, cuBLAS) let frameworks treat the GPU as a fixed target \citep{chen2018tvm,chen2020compilerbasedhardw}.
Inference on phones, laptops, and custom accelerators broke that assumption: different memory structures, different parallelism, and LLM serving costs that made operator-library launch overhead unacceptable \citep{taxbreak2026,kwon2023vllm}.
The compiler's job grew from ``emit a good GPU kernel'' to ``emit a legal, fast schedule for whichever incompatible machine this fleet deploys'' \citep{yirage2026,patel2025llminference}.

\paragraph{Execution models, not speed ranks.}
A GPU oversubscribes warps to hide latency; a CPU hides latency with cache and prefetch; an NPU has no dynamic oversubscription---only a fixed spatial array \citep{nvidia2022hopper,dlbook,amd2024xdna}.
Compilers that pretend an NPU is a slow GPU emit illegal or pathological schedules \citep{wu2024mirage,amd2024aie}.
Table~\ref{tab:hw_landscape} is a taxonomy of models, not a procurement leaderboard \citep{memoryfloor2026}.

\paragraph{Multi-hardware delta.}
The same decoder layer yields a captured fused megakernel on NVIDIA GPU, a cache-blocked ahead-of-time bundle on CPU, and a static tile+DMA ADF on NPU \citep{semianalysis2024tma,amd2024aie,rotem2018glow}.
Identical math, different bytes/token---hardware is a first-class compiler input \citep{yirage2026,chen2020compilerbasedhardw}.

\paragraph{Decode anchors from Chapter~1 (one-line).}
TaxBreak reports hundreds to thousands of launches per token on dense and MoE decoders; removing host overhead bought ${\sim}1.26\times$ on a Qwen-2.5-7B cell without changing math \citep{taxbreak2026,memoryfloor2026,nvidia2024cudagraphs}.
Those numbers are symptoms; this chapter explains the hardware boundaries fusion must respect to move them \citep{patel2025llminference}.

\paragraph{AMD and second-source GPUs.}
AMD CDNA/RDNA paths mirror NVIDIA hierarchy with wavefronts instead of warps and matrix cores reached through HIP \citep{amd2024aie}.
Constraint thinking transfers; magic numbers do not---each SKU row in the hardware model must be populated from datasheets, not copied from H100 \citep{yirage2026,patel2025llminference}.

\paragraph{Latency versus capacity on each class.}
GPU HBM offers enormous bandwidth but only if data stops round-tripping through the hierarchy \citep{nvidia2022hopper,memoryfloor2026}.
CPU DRAM bandwidth is lower but latency to L3 is excellent when blocking is correct \citep{dlbook,patel2025llminference}.
NPU DDR bandwidth is scarce yet deterministic when DMA is planned---predictability can beat peak bandwidth for edge SLAs \citep{amd2024xdna,amd2024aie}.
Decode compilers optimize \emph{where bytes live over time}, not abstract FLOPs \citep{taxbreak2026,kwon2023vllm}.

\paragraph{Operator libraries versus compiler ownership.}
Vendor libraries remain valuable baselines, but LLM decode needs cross-operator residency decisions libraries cannot see \citep{chen2018tvm,taxbreak2026}.
The shift from library calls to compiler-planned megakernels is driven by launches/token and bytes/token, not dislike of cuBLAS \citep{memoryfloor2026,yirage2026}.
Hardware constraints in this chapter explain \emph{why} that shift is mandatory at batch-1 \citep{patel2025llminference}.

\section{Constraint Matrix: The Compiler's Hard Ruler}
\label{sec:constraint_matrix}

% AUTO_VISUAL:tab_constraint_matrix
\begin{table}[t]
\centering
\caption{Cross-hardware constraint matrix.}
\label{tab:constraint_matrix}
\begin{tabular}{lccc}
\toprule
\visTableHeader
\textbf{Dimension} & \textbf{GPU} & \textbf{CPU} & \textbf{NPU} \\
\midrule
On-chip SRAM & Shared mem / L1 & L1/L2/L3 cache & Tile buffers \\
Parallel grain & Warp (32 threads) & Core + SIMD & Spatial PE array \\
Dynamic shapes & Flexible & Flexible & Restricted \\
\bottomrule
\end{tabular}
\end{table}

The compiler does not optimize ``a GPU'' in the abstract---it satisfies a \textbf{constraint matrix}: on-chip capacity, memory hierarchy, parallel grain, shape flexibility \citep{chen2020compilerbasedhardw,yirage2026}.
On-chip SRAM is programmer-managed shared memory on GPU, hardware cache on CPU, explicit tile SRAM on NPU.
Parallel grain is the warp, core+SIMD lane, or spatial processing element (PE) array respectively \citep{nvidia2022hopper,dlbook,amd2024xdna}.

Decode is dominated by \textbf{bandwidth} through the hierarchy: batch-1 re-reads full weights and growing KV each token \citep{kwon2023vllm,taxbreak2026}.
The recurring failure is illegal tier placement---tensors that should stay in fast SRAM spilling to HBM/DDR \citep{memoryfloor2026}.
The matrix lets passes reject illegal fusions \emph{before} codegen instead of in p99 telemetry \citep{patel2025llminference}.

% OPENTIKZ: tri-hardware memory hierarchy
\begin{figure}[t]
\centering
\begin{tikzpicture}[vis base, scale=0.72, transform shape]
 \foreach \x/\title in {-4.8/GPU,0/CPU,4.8/NPU} {
   \node[font=\footnotesize\bfseries] at (\x,2.35) {\title};
 }
 \foreach \x/\lines in {
   -4.8/{Regs\\Shared\\L2\\HBM},
   0/{Regs\\L1/L2/L3\\DRAM},
   4.8/{Tile SRAM\\L2 tile\\DDR}
 } {
   \node[align=center, font=\scriptsize] at (\x,0.8) {\lines};
 }
 \node[font=\scriptsize, visText] at (0,-1.2) {Same math, different fast-memory budgets and scheduling rules};
\end{tikzpicture}
\caption{Tri-hardware memory hierarchy sketch (not to scale): decode planning must read per-class capacities and latencies.}
\label{fig:tri_hw_memory}
\end{figure}

\begin{table}[t]
\centering
\small
\caption{Decode optimization comparison across hardware classes (design-time cheat sheet).}
\label{tab:tri_hw_decode_compare}
\begin{tabular}{@{}p{2.1cm}p{3.2cm}p{3.2cm}p{3.2cm}@{}}
\toprule
\textbf{Dimension} & \textbf{GPU} & \textbf{CPU} & \textbf{NPU} \\
\midrule
Core bottleneck & Launch overhead + HBM bandwidth & Cache misses + NUMA latency & On-chip capacity + static layout \\
Optimization focus & Deep fusion + graph capture & Cache blocking + thread affinity & Static pipeline + DMA plan \\
Best fusion grain & Layer / half-layer stages & Coarse operator groups & Full-layer / multi-layer MegaKernel \\
Dynamic shapes & Strong & Strong & Weak (bucketing/padding) \\
Typical failure & Register/shared spill & LLC eviction + remote NUMA & Compile fail / tile imbalance \\
\bottomrule
\end{tabular}
\end{table}

\paragraph{$\checkmark$ Practice checkpoint.}
Before merging a cross-hardware change, verify the constraint-matrix row for \emph{each} target: capacities, parallel grain, and shape policy \citep{yirage2026}.
A GPU-only green build that breaks NPU legality should fail CI, not ship to heterogeneous fleets \citep{amd2024xdna,patel2025llminference}.

\paragraph{Shape flexibility row.}
GPUs and CPUs handle dynamic sequence lengths and growing KV naturally; static NPUs require bucketing or padding to a finite shape set \citep{amd2024xdna,lai2023relax}.
Relax-style systems made dynamic shapes a first-class compiler concern because they determine \emph{whether} a workload runs, not only how fast \citep{lai2023relax}.
For decode serving, shape policy is as important as SRAM capacity \citep{kwon2023vllm}.

\paragraph{Why ``fastest chip'' is the wrong question.}
Datacenter GPUs lead raw HBM bandwidth but carry kernel launch overhead and cost; CPUs trade bandwidth for control; NPUs lead power efficiency with the least runtime flexibility \citep{taxbreak2026,memoryfloor2026,amd2024aie}.
Each matrix row is a deployment trade---Chapter~24 routes fleets using exactly these cells \citep{patel2025llminference}.

\paragraph{Roofline attachment discipline.}
Attach a roofline or buffer-accounting sheet for each fused region on the target SKU before merge \citep{memoryfloor2026,yirage2026}.
Over-budget regions are build errors: the matrix supplies the numbers; the compiler supplies the inequality check \citep{nvidia2022hopper,chen2020compilerbasedhardw}.

\paragraph{Connecting matrix rows to Part~VI passes.}
Residency passes read on-chip SRAM rows; parallel-grain passes choose warp versus SIMD versus spatial tiling; shape-flexibility passes decide bucketing policies \citep{yirage2026,lai2023relax,chen2020compilerbasedhardw}.
When a pass violates a cell, reject before codegen---the same discipline Chapter~2 asked for on export edges \citep{chen2020compilerbasedhardw,taxbreak2026}.

\paragraph{Using the comparison table in reviews.}
Table~\ref{tab:tri_hw_decode_compare} columns are intentional prompts: if your PR touches GPU fusion, check whether launches/token moved; if it touches CPU blocking, check LLC misses; if it touches NPU ADF, check DMA edges \citep{taxbreak2026,amd2024xdna,dlbook}.
A change that improves one column while regressing another on the same class is incomplete \citep{memoryfloor2026,yirage2026}.
Fleet routers in Part~VII read the same columns when choosing SKU per query class \citep{patel2025llminference}.

\section{GPU Constraints: Capacity and Orchestration Together}
\label{sec:gpu_constraints}

% AUTO_VISUAL:fig_gpu_memory_hierarchy
\begin{figure}[t]
\centering
\begin{tikzpicture}[vis base, scale=0.84, transform shape]
 \node[vis arch node] (s0) at (-5.5,0) {\footnotesize Registers\\256 KB/SM};
 \node[vis arch node] (s1) at (-1.9,0) {\footnotesize Shared mem\\up to 228 KB};
 \node[vis arch node] (s2) at (1.8,0) {\footnotesize L2 cache\\50 MB};
 \node[vis arch node] (s3) at (5.4,0) {\footnotesize HBM\\TB/s class};
 \draw[vis arrow] (s0.east) -- (s1.west);
 \draw[vis arrow] (s1.east) -- (s2.west);
 \draw[vis arrow] (s2.east) -- (s3.west);
\end{tikzpicture}
\caption{GPU memory hierarchy for decode (H100-class sizes).}
\label{fig:gpu_memory_hierarchy}
\end{figure}

\begin{figure}[t]
\centering
\begin{tikzpicture}[vis base, scale=0.88, transform shape]
 \node[vis pipeline node] (rf) at (0,1.2) {\footnotesize Register file};
 \node[vis pipeline node] (smem) at (-2.2,0) {\footnotesize Shared memory};
 \node[vis pipeline node] (tc) at (2.2,0) {\footnotesize Tensor cores (WGMMA)};
 \node[vis arch node] (warp) at (0,-1.1) {\footnotesize Warp scheduler (32-thread grain)};
 \draw[vis arrow] (rf.south) -- (smem.north);
 \draw[vis arrow] (rf.south) -- (tc.north);
 \draw[vis arrow] (warp.north) -- (smem.south);
 \draw[vis arrow] (warp.north) -- (tc.south);
\end{tikzpicture}
\caption{H100 SM resource sketch: fusion budgets contend for the same register and shared-memory pool.}
\label{fig:gpu_sm_resources}
\end{figure}

On an H100 SM, shared memory is configurable up to 228\,KB (227\,KB usable per block after 1\,KB reserved), the register file holds 65{,}536 32-bit registers (256\,KB total, capped at 255 registers per thread), and ${\sim}50$\,MB L2 backs HBM \citep{nvidia2022hopper}.
The \textbf{warp} of 32 threads is the scheduling grain; four warps form a warpgroup for WGMMA; Hopper adds thread-block clusters sharing shared memory across CTAs \citep{shah2024flashattention3}.

\paragraph{Shared memory capacity (constraint loop).}
\textbf{Engineering meaning:} this is the hard ceiling for fused decode tiles---if your fusion region needs 250\,KB, it \emph{will} spill or must split \citep{memoryfloor2026}.
\textbf{Common pitfall:} counting tensor bytes only, ignoring bank padding, double buffering, and epilogue temporaries \citep{nvidia2022hopper}.
\textbf{Debug:} Nsight Compute \texttt{Shared Memory Configuration} and shared-memory throughput; elevated \texttt{local memory traffic} signals register/shared spill to HBM-backed local memory \citep{taxbreak2026}.

\paragraph{Register pressure vs occupancy.}
Resident threads share the 65{,}536-register file; a kernel using $R$ registers per thread caps residents near $65536/R$ \citep{nvidia2022hopper}.
At batch-1, occupancy often matters less than residency---until spills flip the trade \citep{memoryfloor2026}.
\textbf{Debug:} register usage per thread and spill metrics in Nsight Compute; if spills appear, shrink fusion depth or retile before tuning MMA micro-kernels \citep{yirage2026}.

\paragraph{Shared-memory bank conflicts.}
A \textbf{bank conflict} occurs when multiple threads in a warp hit the same shared-memory bank in one cycle, serializing what should be parallel loads \citep{nvidia2022hopper}.
H100 shared memory uses 32 banks of 4 bytes; strided layouts that map many lanes to one bank are a common decode fusion footgun.
Layout is part of residency planning, not a post-pass cosmetic \citep{shah2024flashattention3}.

\paragraph{Kernel launch overhead.}
Each host launch pays dispatch cost; a dozen kernels per layer pays a dozen times per token at batch-1 \citep{taxbreak2026,memoryfloor2026}.
Fusion and CUDA Graph capture attack scheduling overhead directly---but graphs without dataflow still ping-pong inside recorded nodes (Chapter~2) \citep{nvidia2024cudagraphs}.

\paragraph{TMA and async movement (Hopper+).}
\textbf{Engineering meaning:} TMA lets the schedule prefetch tiles into shared memory while math runs on the previous tile---but only if copy edges are static enough to plan \citep{semianalysis2024tma,nvidia2022hopper}.
\textbf{Pitfall:} treating TMA as a band-aid on operator-first ping-pong; bytes/token must still drop \citep{memoryfloor2026}.
\textbf{Debug:} Nsight Systems timelines for TMA overlap; correlate with shared-memory occupancy \citep{semianalysis2024tma}.

\paragraph{Thread-block clusters.}
Hopper clusters let CTAs share shared memory across a larger fusion region---extending Rule~2 when a single block cannot hold the tile \citep{nvidia2022hopper}.
The compiler must prove cluster-wide footprint still fits hardware caps; otherwise cluster use increases complexity without residency wins \citep{yirage2026}.

\paragraph{Multi-hardware delta (GPU vs NPU).}
GPUs hide latency with async copies and deep pipelines; NPUs require those copies to be scheduled statically in the ADF \citep{amd2024xdna,semianalysis2024tma}.
GPU constraint is scheduling overhead \emph{and} capacity; NPU constraint is capacity with zero runtime slack \citep{wu2024mirage}.

\paragraph{WGMMA layout constraint.}
Tensor cores consume operands from shared memory in specific layouts; feeding them from misaligned tiles triggers conversion or gather traffic that negates fusion \citep{shah2024flashattention3,nvidia2022hopper}.
Layout assignment is therefore part of the residency decision, not an afterthought on the critical path \citep{chen2020compilerbasedhardw,yirage2026}.

\paragraph{Chapters~6--8 preview.}
Parts~IV--V tile decode megakernels against the exact numbers in this section---register caps, shared-memory banks, cluster footprints \citep{nvidia2022hopper,semianalysis2024tma}.
When those chapters cite 227\,KB usable shared memory, they mean this inequality, not a metaphor \citep{memoryfloor2026}.

\paragraph{$\checkmark$ Practice checkpoint.}
Before merging GPU fusion: (1)~statically budget shared memory + registers against the target SM; (2)~verify in Nsight Compute no spill and acceptable bank conflicts \citep{nvidia2022hopper,memoryfloor2026}.

\paragraph{Worked example: fused attention tile budget (H100).}
Suppose a decode fusion stage keeps $Q$, $K$, $V$ tiles, softmax scratch, and an output accumulator in shared memory for a head dimension of 128 and block size 64.
A naive byte count might report ${\sim}48$\,KB of tensor data, but the compiler also reserves bank-padding columns, a double buffer for the next $K$ tile prefetch, and a 4\,KB epilogue for the projection matmul fragment \citep{shah2024flashattention3,nvidia2022hopper}.
Add WGMMA operand staging in swizzled layout and the footprint can cross 120\,KB before anyone accounts for cluster-wide sharing \citep{semianalysis2024tma}.
The correct workflow is to build a line-item residency sheet: tensor bytes, padding, pipeline stages, temporaries, then compare the sum to 227\,KB usable \citep{memoryfloor2026,yirage2026}.
If the sheet exceeds the cap, split the fusion at the export edge Chapter~2 marked---for example, keep attention fused but export logits before the MLP rather than forcing an illegal megakernel \citep{chen2020compilerbasedhardw,taxbreak2026}.
Teams that skip the sheet often discover spill only when production batch sizes differ from the lab cell \citep{patel2025llminference}.

\paragraph{Occupancy versus residency trade (when to care).}
High occupancy helps latency hiding when kernels are memory-bound \emph{and} footprints are small \citep{nvidia2022hopper}.
Decode megakernels with large shared-memory tiles intentionally run fewer blocks per SM because the residency win dominates occupancy \citep{memoryfloor2026,dao2022flashattention}.
The mistake is chasing occupancy charts from a GEMM microbench while your fused kernel already spilled \citep{yirage2026}.
Profile bytes/token first; tune occupancy only when spill metrics are clean \citep{taxbreak2026}.

\paragraph{Graph capture interaction.}
CUDA Graphs remove per-launch host overhead but freeze the launch structure \citep{nvidia2024cudagraphs,taxbreak2026}.
Hardware budgets still apply inside captured nodes: a graph does not create more shared memory \citep{nvidia2022hopper,memoryfloor2026}.
Capture after fusion is legal and bytes/token improved---not as a substitute for illegal tiling \citep{yirage2026,chen2020compilerbasedhardw}.
TaxBreak's launch-reduction wins assumed graphs wrapped already-resident schedules \citep{taxbreak2026,memoryfloor2026}.

\paragraph{Blackwell preview (numbers change, inequalities do not).}
Later chapters update register and shared-memory ceilings for Blackwell-class SKUs \citep{nvidia2022hopper,yirage2026}.
The workflow in this section stays identical: line-item budget, compare to modeled cap, reject or split before codegen \citep{chen2020compilerbasedhardw,patel2025llminference}.
New tensor instructions change \emph{how} you fill the budget, not \emph{whether} a budget exists \citep{semianalysis2024tma,shah2024flashattention3}.

\paragraph{L2 as a hidden export edge.}
Even when shared memory is legal, writing full activations through L2 to HBM between subgraphs reintroduces the warehouse trip \citep{memoryfloor2026,kwon2023vllm}.
L2 capacity is large but bandwidth to HBM is still the decode bottleneck at batch-1 \citep{nvidia2022hopper}.
Fusion should minimize round trips, not merely fit in shared memory while exporting every intermediate \citep{chen2020compilerbasedhardw}.

\section{CPU Constraints: Cache, SIMD, and NUMA}
\label{sec:cpu_constraints}

CPU residency is statistical: blocking so tiles fit L1/L2/L3 before eviction, not explicit scratchpad allocation \citep{dlbook,rotem2018glow}.
Parallel grain is core + \textbf{SIMD} (single-instruction multiple-data) lanes---AVX-512 on x86, NEON on Arm---coarser than GPU warps \citep{dlbook}.

\paragraph{NUMA (non-uniform memory access).}
On multi-socket hosts, each socket has local DRAM; remote-socket access adds latency and bandwidth loss \citep{patel2025llminference}.
\textbf{Engineering meaning:} pin decode threads and weight buffers to the same NUMA node, or replicate read-only weights per node when cheaper than remote traffic \citep{memoryfloor2026}.
\textbf{Debug:} \texttt{numastat}, \texttt{perf} LLC miss counters, and cross-node traffic on the decode loop.

\paragraph{Bandwidth floor.}
A 7B model in BF16 moves ${\sim}14$\,GB of weights per token at batch-1---dividing by DRAM bandwidth sets a latency floor no micro-kernel beats \citep{kwon2023vllm,memoryfloor2026}.
CPU wins on flexibility and cost at low QPS, not on matching H100 tail latency \citep{patel2025llminference}.

\paragraph{SIMD + quantization coupling.}
Int8 weights shrink cache footprint and fill SIMD lanes denser---double win when bandwidth-bound \citep{dlbook}.
GPU fusion depth copied blindly can evict KV working sets; CPU schedules need cache-aware blocking \citep{rotem2018glow}.

\paragraph{$\checkmark$ Practice checkpoint.}
Validate blocking plans with measured L3 miss rate and bytes/token at the production batch policy---paper budgets are insufficient on CPUs \citep{memoryfloor2026,dlbook}.

\paragraph{Prefetch and layout.}
Sequential weight streams benefit from hardware prefetch when layouts are contiguous \citep{dlbook}.
Fragmented layouts defeat prefetch and amplify bandwidth tax on a already floor-limited CPU decode cell \citep{patel2025llminference,kwon2023vllm}.

\paragraph{Multi-hardware delta (CPU vs GPU).}
CPU residency is statistical---miss rates matter; GPU residency is contractual---kilobytes over cap spill \citep{nvidia2022hopper,dlbook}.
A blocking plan that ``should fit'' in L3 must be validated; a shared-memory plan that exceeds 227\,KB must be rejected arithmetically \citep{memoryfloor2026,yirage2026}.

\paragraph{When CPU decode makes business sense.}
Low-QPS endpoints, strict cost caps, and workloads with heavy control flow may favor CPU even when GPUs lead peak throughput \citep{patel2025llminference,kwon2023vllm}.
The constraint chapter equips you to recognize those regimes instead of forcing GPU schedules everywhere \citep{rotem2018glow,dlbook}.

\paragraph{Glow and AOT mental model.}
Ahead-of-time bundles with static shapes mirror NPU thinking on CPU: fewer surprises at runtime, more work at compile time \citep{rotem2018glow,chen2020compilerbasedhardw}.
The book's CPU path favors predictable residency over last-minute operator dispatch \citep{yirage2026}.

\paragraph{Worked example: L3 blocking for a 7B MLP slice.}
At batch-1, a single MLP projection reads the full weight matrix from DRAM each token \citep{kwon2023vllm}.
Blocking chooses a strip height so the active weight tile plus input vector plus output accumulator stays in L3 across the inner dimension \citep{dlbook,rotem2018glow}.
If the strip is too wide, the next layer's KV working set evicts the weight tile---exactly the 15\% regression seen when GPU fusion depth was copied without cache accounting \citep{patel2025llminference}.
Measure LLC misses per token at the production thread count; if misses scale with layer count, the block is too aggressive \citep{memoryfloor2026}.
CPU optimization is iterative measurement against a byte budget, not a one-shot formula \citep{dlbook}.

\paragraph{Thread pinning recipe.}
Bind decode worker threads to cores on the same NUMA node as the weight buffer \citep{patel2025llminference}.
On dual-socket hosts, replicate read-only weights per socket when replication cost is less than remote DRAM traffic during the decode loop \citep{memoryfloor2026}.
Document the pinning policy beside the blocking plan---silent migration to a remote socket looks like ``compiler regression'' in dashboards \citep{yirage2026}.

\paragraph{AVX-512 and decode tails.}
Wide SIMD helps when inner dimensions are multiples of the vector width; odd head sizes leave lanes idle unless the compiler pads or peels tails \citep{dlbook}.
Padding changes cache footprint; peeling adds scalar epilogue cost \citep{rotem2018glow}.
Co-decide head padding policy with quantization and blocking, not in a late vectorization pass \citep{chen2020compilerbasedhardw}.

\section{NPU Constraints: Static Spatial Plans}
\label{sec:npu_constraints}

Spatial NPUs are deterministic and unforgiving.
AMD XDNA2/AIE-ML tiles expose 64\,KB L1 SRAM per compute tile (eight 8\,KB banks) with DMA, column L2 tiles at 512\,KB, and DDR via shim tiles---no hardware cache \citep{amd2024xdna,amd2024aie}.
What is resident is exactly what the compile-time DMA schedule placed.

\paragraph{Tile capacity (constraint loop).}
\textbf{Engineering meaning:} operands must fit tile SRAM or be split with explicit multi-tile DMA chains---there is no runtime cache to hide mistakes \citep{amd2024xdna}.
\textbf{Failure mode:} compile failure or severely imbalanced pipelines that idle tiles \citep{wu2024mirage}.
\textbf{Debug:} compiler legality reports and DMA edge counts per token; compare against modeled tile budgets in YiRage \citep{yirage2026}.

\paragraph{Spatial cascade.}
Adjacent tiles pass data through cascade links without DDR round trips---the on-chip analogue of deep software pipelining \citep{amd2024aie}.
The spatial plan is fixed at compile time; dynamic decode loops require shape bucketing or padding \citep{lai2023relax,amd2024xdna}.

\paragraph{$\checkmark$ Practice checkpoint.}
NPU promotions require green static compile, balanced tile utilization, and bytes/token regression versus the prior ADF hash---no ``works on GPU'' excuses \citep{yirage2026,amd2024xdna}.

\paragraph{Optimization means balance, not search.}
On GPUs, imperfect schedules hide behind dynamic scheduling; on NPUs, imbalanced tile assignment idles silicon with no runtime recovery \citep{wu2024mirage,amd2024xdna}.
Edge deployments buy deterministic latency and power at the price of compile-time rigor \citep{amd2024aie,rotem2018glow}.

\paragraph{Attention on spatial fabrics.}
Canonical edge fusion maps attention as a spatial sweep with scores resident in tile SRAM and DDR touched only for inputs/outputs---the spatial analogue of a GPU megakernel \citep{dao2022flashattention,memoryfloor2026,amd2024xdna}.

\paragraph{DDR is not ``slow GPU HBM.''}
Off-chip DRAM bandwidth on edge cards is scarce; every DMA must be planned, counted, and minimized in the ADF \citep{amd2024xdna,amd2024aie}.
Bytes/token ledgers from Chapter~1 apply unchanged---only the mechanism counter becomes DMA edges instead of launches \citep{taxbreak2026,memoryfloor2026}.

\paragraph{Mirage and spatial compilers.}
Spatial compiler stacks (Mirage-class) search tile assignments where GPU stacks search thread blocks \citep{wu2024mirage,amd2024xdna}.
The mindset shift from Chapter~2 is the same; the search space geometry changes \citep{chen2020compilerbasedhardw,yirage2026}.

\paragraph{ADF planning checklist (edge NPU).}
Start from a shape bucket the product will ship---for example ctx=2048 with fixed head count \citep{lai2023relax,amd2024xdna}.
Assign each tensor to a tile SRAM bank with explicit DMA in and out edges; count bytes moved per token \citep{amd2024aie}.
Balance compute across tiles: an attention sweep that leaves half the array idle fails even when compile succeeds \citep{wu2024mirage}.
Run legality first, then bytes/token, then ms/token---dynamic GPU shortcuts do not apply \citep{yirage2026,memoryfloor2026}.
Version the ADF hash in CI next to GPU kernel hashes so regressions are comparable across classes \citep{patel2025llminference}.

\paragraph{When NPU compile fails (readable errors).}
Out-of-tile SRAM is the most common failure: the fix is smaller spatial tiles or fewer fused stages, not ``retry with more threads'' \citep{amd2024xdna}.
DMA stride misalignment fails legalization because the shim expects burst-friendly layouts \citep{amd2024aie}.
Treat compile errors as constraint violations with the same seriousness as shared-memory overflow on GPU \citep{chen2020compilerbasedhardw,yirage2026}.

\section{Hardware-Aware Compilation: Constraints Become Pass Inputs}
\label{sec:hardware_aware_compile}

Hardware-aware compilation makes constraints data, not folklore.
Three coupled decisions dominate: \textbf{tiling}, \textbf{layout}, and \textbf{pass} ordering \citep{chen2020compilerbasedhardw,lai2023relax}.

\textbf{Tiling} sizes on-chip working sets to fit 228\,KB shared memory (GPU), a cache level (CPU), or 64\,KB tile SRAM (NPU)---the largest bytes/token lever \citep{nvidia2022hopper,amd2024aie}.
\textbf{Layout} ensures coalesced global access, bank-conflict-free shared tiles, and DMA-friendly strides \citep{shah2024flashattention3}.
\textbf{Pass order} matters: layout-after-fusion transposes can re-materialize activations fusion meant to keep resident \citep{lai2023relax,chen2018tvm}.
Precision (BF16/FP8/int8) co-decides with tiling because formats change footprint and native instruction support per SKU \citep{shah2024flashattention3,liu2024deepseekv3}.

\begin{figure}[t]
\centering
\begin{tikzpicture}[vis base, scale=0.78, transform shape]
 \node[vis pipeline node] (m) at (-4.2,0) {\footnotesize Hardware model\\(capacities + costs)};
 \node[vis arch node] (t) at (-1.4,0) {\footnotesize Tiling / fusion};
 \node[vis arch node] (l) at (1.4,0) {\footnotesize Layout / vectorize};
 \node[vis pipeline node] (v) at (4.2,0) {\footnotesize Legality check\\+ codegen};
 \draw[vis arrow] (m) -- (t) -- (l) -- (v);
 \node[font=\scriptsize, visText] at (0,-1.0) {Residency budget informs all stages jointly};
\end{tikzpicture}
\caption{Hardware-aware compile flow: the chip model feeds tiling, layout, and pre-codegen legality checks.}
\label{fig:hw_aware_compile_flow}
\end{figure}

\paragraph{$\checkmark$ Practice checkpoint.}
Diff IR or assembly for new global buffers; count launches/token (GPU) and DMA edges (NPU); confirm tiling still fits modeled capacity per target \citep{taxbreak2026,yirage2026,memoryfloor2026}.

\paragraph{Joint residency budget.}
Larger tiles improve compute intensity but consume SRAM; padded layouts avoid bank conflicts but cost bytes; double buffering for pipelines doubles footprints \citep{nvidia2022hopper,shah2024flashattention3}.
Hardware-aware compilation is one joint optimization against modeled capacity---not tiling, then layout, then surprise spill \citep{chen2020compilerbasedhardw,yirage2026}.
Passes that fix tiling first and discover layout incompatibility later must backtrack; residency budgets should constrain all three from the start \citep{lai2023relax}.

\paragraph{Per-target pass forks.}
The right pass order differs by class: NPUs need tile assignment before vectorization details; GPUs may fuse before layout on some graphs but not others \citep{chen2018tvm,chen2020compilerbasedhardw}.
Expose pipelines explicitly (Part~VI) so hardware constraints become visible engineering choices, not accidental emergent behavior \citep{yirage2026}.

\paragraph{Anti-pattern: layout-after-fusion transpose.}
A transpose inserted after fusion may force a full activation write/read pair---exactly the export Chapter~2 removed \citep{lai2023relax,chen2020compilerbasedhardw,memoryfloor2026}.
Hardware-aware compilers order layout and fusion so residency commitments remain satisfiable \citep{yirage2026,taxbreak2026}.

\paragraph{Anti-pattern: precision change without retile.}
Dropping BF16 activations to FP8 halves bytes but may change vectorization and shared-memory padding requirements \citep{shah2024flashattention3,liu2024deepseekv3}.
Co-decide precision with tiling when bandwidth-bound \citep{memoryfloor2026,patel2025llminference}.

\section{Common Cross-Hardware Migration Mistakes}
\label{sec:migration_mistakes}

Teams porting decode schedules between silicon classes repeat the same failures.
This section names five patterns so you can catch them in design review before they become production incidents \citep{patel2025llminference,yirage2026,chen2020compilerbasedhardw}.

\paragraph{Mistake 1: capacity math from the wrong machine.}
Using H100 shared-memory numbers on a laptop NPU compile, or assuming CPU L3 is as programmable as GPU shared memory, produces schedules that are either illegal or silently slow \citep{nvidia2022hopper,amd2024xdna,dlbook}.
\textbf{Fix:} read the constraint-matrix row for \emph{each} target before writing fusion rules \citep{yirage2026}.

\paragraph{Mistake 2: promoting GPU fusion depth to CPU unchanged.}
Deep GPU megakernels maximize on-chip residency when spill is impossible; on CPU the same depth evicts KV and weight tiles from LLC \citep{rotem2018glow,patel2025llminference,memoryfloor2026}.
\textbf{Fix:} re-block for cache hierarchy; accept shallower fusion with better miss rates \citep{dlbook}.

\paragraph{Mistake 3: expecting dynamic shapes on static NPUs.}
Serving growth in sequence length without re-bucketing breaks ADF plans that assumed fixed ctx \citep{lai2023relax,amd2024xdna}.
\textbf{Fix:} maintain a finite shape set; compile one ADF per bucket; pad explicitly where product allows \citep{yirage2026}.

\paragraph{Mistake 4: optimizing peak FLOPs after fusion is legal.}
A legal kernel can still lose to scheduling overhead if launches/token remain high, or lose to bandwidth if bytes/token did not drop \citep{taxbreak2026,memoryfloor2026}.
\textbf{Fix:} gate promotions on Chapter~1 counters, not vendor GEMM banners \citep{patel2025llminference}.

\paragraph{Mistake 5: hard-coded magic numbers in kernels.}
Scattering 228\,KB constants in CUDA blocks Blackwell retargeting and breaks CI legality checks \citep{yirage2026,iree2024,nvidia2022hopper}.
\textbf{Fix:} centralize capacities in the hardware model; let passes read limits \citep{chen2020compilerbasedhardw}.

\paragraph{Design-review prompt.}
Ask in every cross-hardware PR: \emph{Which constraint cell changed, what is the new byte budget, and which mechanism counter proves the win?} \citep{taxbreak2026,memoryfloor2026,patel2025llminference}
If the author cannot answer, the change is not ready to merge \citep{yirage2026}.

\paragraph{Postmortem template.}
When production latency regresses after a fusion promotion, capture: SKU, ctx, batch, static byte budget, measured shared-memory or tile usage, spill or DMA flags, and before/after bytes/token \citep{memoryfloor2026,taxbreak2026,nvidia2022hopper}.
Most postmortems reveal a violated inequality from this chapter, not a mysterious hardware flake \citep{patel2025llminference,yirage2026}.
Store postmortems beside benchmark worksheets so the next engineer does not repeat the mistake \citep{chen2020compilerbasedhardw}.

\section{Cross-Hardware Benchmarking: Fair Comparison}
\label{sec:benchmark_method}

% AUTO_VISUAL:fig_benchmark_method
\begin{figure}[t]
\centering
\begin{tikzpicture}[vis base, scale=0.9, transform shape]
 \node[vis pipeline node] (s0) at (-4.0,0) {\footnotesize Pin model+ctx\\BS${=}1$};
 \node[vis arch node] (s1) at (0.0,0) {\footnotesize Normalize\\bytes/token, launches/token};
 \node[vis pipeline node] (s2) at (4.0,0) {\footnotesize Compare GPU/CPU/NPU\\same metrics};
 \draw[vis arrow] (s0.east) -- (s1.west);
 \draw[vis arrow] (s1.east) -- (s2.west);
\end{tikzpicture}
\caption{Cross-hardware benchmark flow used in this book.}
\label{fig:benchmark_method}
\end{figure}

Honest hardware comparison fixes the workload, then grades decode metrics that predict serving cost---not vendor peak FLOPs \citep{taxbreak2026,memoryfloor2026,patel2025llminference}.
\textbf{launches/token} isolates scheduling overhead; \textbf{bytes/token} isolates bandwidth/residency \citep{kwon2023vllm}.
Decomposed counters are actionable; a fast GEMM microbench is not \citep{memoryfloor2026}.

\paragraph{Per-class mechanisms.}
GPUs move launches/token with fusion and graphs; CPUs move LLC misses and NUMA traffic; NPUs move DMA edge counts \citep{nvidia2022hopper,amd2024aie,dlbook}.
Report the \emph{same} headline metrics while diagnosing class-specific mechanisms underneath \citep{taxbreak2026}.

\paragraph{Reproducibility.}
Pin weights, framework versions, context buckets, and log configs so SKU A compares to SKU B and to next month's compiler \citep{memoryfloor2026,yirage2026}.
Appendix~B expands long-context and fleet-batch variants; this section states the core decode-cell discipline \citep{kwon2023vllm}.

\paragraph{Decomposition beats single scores.}
End-to-end ms/token alone does not say whether to fuse or quantize; launches/token points to scheduling work; bytes/token points to residency work \citep{taxbreak2026,memoryfloor2026}.
This decomposition is why Chapter~1 insisted on goodput counters alongside latency \citep{patel2025llminference,kwon2023vllm}.

\paragraph{Autobooks harness analogy.}
The book's own loop fixes a harness, scores every change on the same metrics, and reverts regressions---benchmark discipline as engineering culture, not a one-off table \citep{yirage2026,taxbreak2026}.
Cross-vendor comparisons without pinned methodology decay into marketing slides \citep{memoryfloor2026}.

\paragraph{Fleet vs batch-1 cells.}
Report batch-1 mindset counters separately from fleet goodput---continuous batching changes weight reuse but not per-request export edges inside layers \citep{kwon2023vllm,patel2025llminference}.
A schedule that helps interactive tail latency but fails at fleet averages still needs framework scheduling work \citep{memoryfloor2026,taxbreak2026}.

\paragraph{Long-context buckets.}
Benchmark ctx $\in\{512,2048,8192\}$ minimum; paging can change effective residency at 32K even when 2K looked healthy \citep{kwon2023vllm,dao2022flashattention}.
Hardware-aware compilers must re-check tier budgets when block tables grow \citep{yirage2026,memoryfloor2026}.

\paragraph{$\checkmark$ Practice checkpoint.}
Ship benchmark artifacts with every hardware-target promotion: bytes/token, launches/token, ms/token at ctx $\in\{512,2048,8192\}$, plus the mechanism counter for that class \citep{taxbreak2026,memoryfloor2026}.

\paragraph{Benchmark worksheet (copy into CI).}
Record model id, weight format, framework build, SKU, driver, ctx, batch policy, and power cap \citep{memoryfloor2026,yirage2026}.
Run warm-up tokens, then collect median ms/token over at least 100 tokens at batch-1 for interactive cells \citep{taxbreak2026}.
Log launches/token from the profiler trace; log bytes/token from a residency ledger or instrumented allocator \citep{kwon2023vllm,memoryfloor2026}.
Attach the worksheet to the PR; reviewers compare mechanism counters, not adjectives \citep{patel2025llminference}.
When a GPU win does not move bytes/token, suspect measurement noise or hidden exports \citep{chen2020compilerbasedhardw}.

\paragraph{Sanity checks before claiming cross-vendor victory.}
Same tokenizer, same weight checkpoint, same numerical policy (BF16 vs FP8), and same ctx buckets \citep{liu2024deepseekv3,memoryfloor2026}.
Disable unrelated framework features (speculative decoding, paging) unless the product ships them \citep{kwon2023vllm}.
Report confidence intervals when variance exceeds 2\%---decode cells are noisy at small batch \citep{patel2025llminference}.

\section{YiRage Hardware Modeling: One Pass Stack, Many SKUs}
\label{sec:yirage_modeling}

\textbf{YiRage} encodes this chapter as per-target \textbf{chip architecture models}---capacities, parallelism grains, latency/bandwidth ratios, launch floors---that residency, pipeline, and sync passes read and check \citep{yirage2026,chen2020compilerbasedhardw}.
A fusion's footprint is compared to modeled tier capacity; tile shapes are validated against SRAM and bank structure; launch and scheduling costs inform schedule ranking \citep{wu2024mirage}.

Updating a model re-specializes passes when a new SKU ships---H100 $\rightarrow$ Blackwell register growth, for example---without rewriting every hand kernel \citep{iree2024,patel2025llminference}.
Legality-before-codegen is the scalability property heterogeneous fleets need: illegal residency fails compile instead of silently spilling in production \citep{yirage2026,amd2024xdna}.

\paragraph{$\checkmark$ Practice checkpoint.}
Treat the hardware model as versioned data: diff model entries in CI when SKUs change; reject passes that hard-code capacities instead of reading the model \citep{yirage2026}.

\paragraph{Capacities plus costs.}
Models need latencies, bandwidths, launch floors, and tensor-core throughput ratios---not SRAM kilobytes alone---to \emph{rank} legal schedules, not only accept/reject them \citep{nvidia2022hopper,shah2024flashattention3,wu2024mirage}.
A legal but slow schedule is still a product failure \citep{patel2025llminference,memoryfloor2026}.

\paragraph{Portability across SKUs.}
Because passes read the model, Blackwell register growth or new tensor instructions become data updates rather than rewriting every kernel in the tree \citep{yirage2026,iree2024}.
Hardware knowledge lives in one auditable place instead of rotting inside hand-tuned CUDA \citep{chen2020compilerbasedhardw}.

\paragraph{IREE and modular retargeting.}
IREE-style retargeting treats backends as plugins over shared IR---the same architectural pattern YiRage uses for decode fusion \citep{iree2024,yirage2026,chen2020compilerbasedhardw}.
Hardware models are the plugin contract \citep{patel2025llminference}.

\paragraph{Chapters~6--9 dependency.}
Residency, role assignment, pipeline formation, and synchronization passes in Part~IV all consume the chip model introduced here \citep{yirage2026,wu2024mirage}.
Read this chapter once carefully; later chapters reference these inequalities by label rather than re-deriving them \citep{nvidia2022hopper,amd2024xdna}.

\paragraph{Model schema (conceptual).}
Each SKU entry typically carries tier capacities (registers, shared memory or tile SRAM, L2, DRAM), parallelism grains (warp size, SIMD width, tile array), bandwidth/latency ratios, launch or DMA floor costs, and supported instruction sets (WGMMA, matrix cores, int8 dot) \citep{yirage2026,nvidia2022hopper,amd2024xdna}.
Passes query the schema; kernels never embed literals \citep{chen2020compilerbasedhardw}.
When Blackwell raises register file limits, bump the model version and re-run legality on the fusion corpus---no manual kernel sweep \citep{iree2024,patel2025llminference}.

\paragraph{Fleet heterogeneity.}
Production may run Hopper datacenter GPUs, CPU fallback nodes, and edge NPUs simultaneously \citep{patel2025llminference,amd2024aie}.
One IR, multiple chip models, divergent legal schedules---that is normal, not a failure of the compiler architecture \citep{yirage2026,wu2024mirage}.
Chapter~24 returns to routing policies; this chapter supplies the per-target numbers those policies consume \citep{memoryfloor2026}.

\paragraph{Testing hardware models.}
Unit-test the model layer: given a synthetic fusion footprint, assert the legality pass accepts or rejects against golden SKU rows \citep{yirage2026,chen2020compilerbasedhardw}.
Regression tests caught more fleet incidents than microbench suites because they encode inequalities directly \citep{patel2025llminference,memoryfloor2026}.
When vendors ship stepping changes, update golden rows and re-run the corpus before merging driver pins \citep{iree2024,nvidia2022hopper}.

\section{Key Takeaways}
\label{sec:ch03_key_takeaways}

\begin{itemize}
\item \textbf{Quantify constraints, never hand-wave.}
``Fast GPU memory'' is not actionable; 227\,KB usable shared memory per H100 block is \citep{nvidia2022hopper}.
\textbf{Action:} attach static byte budgets to every fusion PR; spill in Nsight Compute is a build failure, not a tuning opportunity \citep{memoryfloor2026}.

\item \textbf{No universal optimal schedule.}
GPU layer fusion, CPU cache blocking, and NPU static MegaKernels are different artifacts for the same math \citep{amd2024xdna,rotem2018glow,mpk2025}.
\textbf{Action:} fork compile pipelines per target; never promote GPU-only evidence to fleet-wide schedules \citep{patel2025llminference,yirage2026}.

\item \textbf{Decode optimizes memory first, FLOPs second.}
Batch-1 limits are bytes and launches on every class \citep{taxbreak2026,kwon2023vllm,memoryfloor2026}.
\textbf{Action:} rank backlog items by bytes/token and launches/token before MMA tile sweeps \citep{patel2025llminference}.

\item \textbf{Hardware belongs in a configurable model.}
Capacities and costs as data let passes specialize and check legality \citep{yirage2026,wu2024mirage}.
\textbf{Action:} add new SKUs by extending the chip model, not scattering magic numbers in kernels \citep{iree2024}.

\item \textbf{Benchmark with decode metrics, not peak FLOPs.}
Normalize bytes/token and launches/token across GPU, CPU, and NPU cells \citep{taxbreak2026,memoryfloor2026}.
\textbf{Action:} store paired worksheets in CI; reject promotions that move only microbench scores \citep{patel2025llminference}.
\end{itemize}

\section{Conclusion}
\label{sec:ch03_conclusion}

This chapter grounded Chapter~2's dataflow mindset in numbers: GPU shared memory and registers with launch overhead, CPU caches and NUMA, NPU tile SRAM and static spatial dataflow \citep{nvidia2022hopper,dlbook,amd2024xdna,amd2024aie}.
Decode is bandwidth- and scheduling-bound on all three; hardware-aware tiling, layout, and pass order---guided by explicit models---turn diagnosis into shippable kernels \citep{chen2020compilerbasedhardw,yirage2026,taxbreak2026,memoryfloor2026}.

With the cross-class frame set, Chapter~\ref{chap:ch04} descends into NVIDIA Hopper/Blackwell CUDA: registers, TMA, thread-block clusters, and tensor cores as concrete optimization bounds for decode megakernels \citep{nvidia2022hopper,semianalysis2024tma}.
Bring your constraint-matrix row and export-edge worksheet from Chapters~1--2; Chapter~4 shows how to spend those bytes on silicon \citep{memoryfloor2026}.

\paragraph{Worksheet gate.}
Before opening Hopper-specific chapters, fill Table~\ref{tab:tri_hw_decode_compare} for your deployment targets and note which bottleneck column dominates \citep{taxbreak2026,memoryfloor2026,patel2025llminference}.
If you cannot name launches/token and bytes/token for your SKU, return to Chapter~1 before studying TMA details \citep{nvidia2022hopper,semianalysis2024tma}.

\paragraph{Hands-on next step.}
Pick one fused region in your decoder and compute shared-memory + register bytes statically; compare to the H100 row in Figure~\ref{fig:gpu_memory_hierarchy} \citep{nvidia2022hopper}.
If over budget, split stages per Chapter~2's five-step playbook before tuning MMA configs \citep{chen2020compilerbasedhardw,yirage2026,memoryfloor2026}.

\paragraph{Onboarding sentence for teams.}
Post in your compiler design doc: \emph{``We do not merge fusion without a hardware-model legality check and decode-cell counters.''} \citep{yirage2026,taxbreak2026,memoryfloor2026}
Culture follows the checklist in CI, not the poster alone \citep{patel2025llminference}.

\paragraph{Part~I closure.}
Chapters~1--3 formed the measurement, mindset, and constraint foundation for everything that follows \citep{taxbreak2026,memoryfloor2026,chen2020compilerbasedhardw}.
You can now name what to optimize (bytes and launches), how to think (residency before exports), and what must fit (per-class capacities) \citep{yirage2026,patel2025llminference}.
Part~II begins the GPU descent: same inequalities, more silicon detail, more concrete pass examples \citep{nvidia2022hopper,semianalysis2024tma}.

\paragraph{Glossary carry-forward.}
Terms introduced here recur throughout the book: \textbf{SIMT}, \textbf{warp}, \textbf{WGMMA}, \textbf{TMA}, \textbf{NUMA}, \textbf{DMA}, and \textbf{spatial dataflow} \citep{nvidia2022hopper,amd2024xdna,dlbook}.
When later chapters use them without redefinition, refer back to this chapter's first-use explanations \citep{yirage2026}.
Consistent vocabulary across GPU, CPU, and NPU sections prevents teams from talking past each other in design reviews \citep{patel2025llminference,chen2020compilerbasedhardw}.

\typeout{END_CHAPTER "ch03" \theabspage}
